\href{https://github.com/yassienE4/OS-VM/commit/929f4ec3803c7845995459e7b9b7fca693f79cf8}{\ul{https://github.com/yassienE4/OS-VM/commit/929f4ec3803c7845995459e7b9b7fca693f79cf8}}

How to run in command mode

\begin{enumerate}
\def\labelenumi{\arabic{enumi}.}
\item
  cd The-Shell
\item
  ./run.sh main.cpp
\end{enumerate}

How to run in file mode

\begin{enumerate}
\def\labelenumi{\arabic{enumi}.}
\item
  cd The-Shell
\item
  g++ -std=c++17 -w main.cpp -o main
\item
  ./main args
\end{enumerate}

The code takes a txt file as args which should include commands

Commands.txt is included as an example

The code works by first checking the arguments that the code was started
with, if no arguments it starts a loop to get commands via inputs, while
if there are args it looks for a file and reads the commands line by
line from that file.

It then parses the commands into a custom data structure i created,
commands.h

string name;

vector\textless string\textgreater{} args;

string inFile;

string outFile;

bool append;

bool background;

bool isBuildIn;

The data structure allows me to understand the data and execute the
correct command according to if there is a built in version, it runs in
the background, it reads or writes to a file etc.

The commands are then executed either via forking and then executing or,
by implementing the behaviors using C system functions
